% Generated draft: requires mathematical and source review.
\documentclass[11pt]{article}
\usepackage[margin=25mm]{geometry}
\usepackage{fontspec}
\setmainfont{Latin Modern Roman}
\usepackage{amsmath,amssymb,mathtools,unicode-math}
\setmathfont{Latin Modern Math}
\usepackage{enumitem,needspace,xurl,hyperref}
\setlength{\parindent}{0pt}
\setlength{\parskip}{5pt}
\setlength{\emergencystretch}{3em}
\newcommand{\N}{\mathbb{N}}
\newcommand{\Z}{\mathbb{Z}}
\newcommand{\Q}{\mathbb{Q}}
\newcommand{\R}{\mathbb{R}}
\newcommand{\C}{\mathbb{C}}
\newcommand{\F}{\mathbb{F}}
\newcommand{\A}{\mathbb{A}}
\newcommand{\PP}{\mathbb{P}}
\newcommand{\E}{\mathbb{E}}
\newcommand{\eps}{\varepsilon}
\newcommand{\dd}{\,\mathrm{d}}
\begin{document}
\section*{UnsolvedMath ID: 30001116}
\textbf{Source problem code:} catalogue entry
\subsection{Definitions and mathematical statement}
Let \(M\) denote a disordered pinning or renewal model of the kind intended by the source, and let \(\beta\) be an admissible disorder parameter for that model. The source does not specify the class of models or the range of \(\beta\). Write
\[
 h_c(M,\beta):=h_c(\beta),\qquad
 h_c^{\mathrm a}(M,\beta):=h_c^{\mathrm a}(\beta)
\]
for its quenched and annealed critical points, respectively. The original statement writes the annealed quantity as \(h_{ac}(\beta)\). Quenched quantities describe a fixed disorder realization before any averaging over disorder, whereas annealed quantities incorporate disorder averaging in the model's statistical weights. The partition function, disorder distribution, averaging procedure, and critical-point convention are not supplied, so these descriptions do not determine particular critical points.

Where both critical points are finite, define their difference by
\[
 \Delta_M(\beta):=h_c(M,\beta)-h_c^{\mathrm a}(M,\beta).
\]
Let \(\alpha_M\) denote the renewal-law tail exponent, if one is specified for \(M\). No convention for this exponent, tail asymptotic, or admissible range is supplied. Let \(\nu_{\mathrm q}(M,\beta)\) and \(\nu_{\mathrm a}(M,\beta)\) stand for corresponding quenched and annealed critical exponents, whenever the observable and limiting procedure defining them have been fixed and the exponents exist. These symbols are placeholders: the source does not identify the critical observable or define its exponent.

\textbf{Source question.} Determine, for the intended models and admissible disorder parameters, whether the quenched and annealed critical points differ and whether their corresponding critical exponents differ. Where the displayed quantities are defined, the comparisons are
\[
 \Delta_M(\beta)\stackrel{?}{=}0,
 \qquad
 \nu_{\mathrm q}(M,\beta)\stackrel{?}{=}
 \nu_{\mathrm a}(M,\beta).
\]
How do these comparisons depend on the renewal-law tail exponent \(\alpha_M\)?

The source poses a broad classification question, not a specified universal equality or inequality. A precise model-specific statement requires the missing model, disorder, renewal-tail, parameter-domain, and critical-exponent definitions. No boundary or marginal-tail cases are specified, and none are excluded here.
\subsection{Short English statement}
For disordered pinning or renewal models, compare the quenched and annealed critical points and corresponding critical exponents, and determine how their agreement or disagreement depends on the renewal-law tail exponent. The supplied record leaves the model class, disorder assumptions, parameter ranges, tail convention, and meaning of critical exponent unspecified; this is therefore a qualified statement of the source question requiring mathematical review.
\subsection{Sources}
\begin{itemize}
\item Supplied local UnsolvedMath registry record, permanent ID \(30001116\), source problem code \texttt{catalogue entry}, titled ``Disorder Relevance in Pinning and Renewal Models,'' in the Oberwolfach Reports Open Problems collection. The record gives proposed year 2008 and names no proposer.
\item Supplied source citation: \emph{Interplay of Analysis and Probability in Physics} (2009), Oberwolfach Reports, pp. 3065--3138, DOI \texttt{10.4171/owr/2008/55}. Original OWR report: \url{https://doi.org/10.4171/owr/2008/55}. The supplied background describes these reports as records of mathematical workshops and their open-problem discussions.
\item Supplied literature source, labeled ``finite-range correlated pinning'': \url{https://doi.org/10.1016/j.spa.2012.06.007}.
\item Supplied literature source, labeled ``quenched and annealed pinning curves'': \url{https://arxiv.org/abs/0903.3704}.
\item Literature assessment checked on \textbf{2026-08-21}: this is supplied dated curation, not independently verified here. The record labels the problem \texttt{partially\_solved} and reports that quenched and annealed critical points and exponents coincide in disorder-irrelevant regimes and differ in relevant regimes for several pinning models, while no universal answer to the broad question is available. The background characterizes this as dated literature triage derived from the supplied OWR curation and requests independent verification of time-sensitive status and cited claims.
\item The registry labels the statement \texttt{corrected\_verified} and reports that its mathematical content was checked against the official Oberwolfach report, with notation, context, and question form normalized without changing the source-level claim. This is a supplied verification assertion, not an independent verification performed for this draft.
\end{itemize}
\end{document}
